\documentclass[twocolumn]{aastex631}
\begin{document}
\title{The reionization ionizing-photon-budget shortfall for star-forming galaxies is not robust to systematics at z\~{}7 (beta systematic corner)}
\author{NebulaMind Lab (autonomous pipeline)}
\affiliation{NebulaMind Lab, \url{https://lab.nebulamind.net}}
\begin{abstract}
Reionization ionizing-photon-budget reconciliation at z\~{}7: star-forming galaxies require f\_esc=0.105 (+0.106/-0.054) to close the budget (Madau-Dickinson SFRD, log xi\_ion=25.5+/-0.15, clumping C=2-5, JWST-SFRD tail), versus indirect-proxy-inferred f\_esc=0.050 (+0.076/-0.030) from LzLCS O32/beta calibrations. Median delta(required-inferred)=+0.048 dex-frac (16-84\%: -0.034 to +0.155); 73\% of the systematic MC shows a shortfall. Conclusion: the budget CLOSES within the systematic, and the sign holds under both O32 and beta calibrations. The result is bounded by the xi\_ion x clumping x proxy-calibration systematic, not by statistics. Generated autonomously from public data (jwst) via the NebulaMind Lab runner: a bounded, reproducible, descriptive study.
\end{abstract}
\section{Introduction}
Introduction:
Recent studies have highlighted a potential crisis in the reionization photon budget, suggesting that star-forming galaxies may not produce enough ionizing photons to account for the observed reionization process [Muñoz2024]. This discrepancy has sparked interest in reconciling the ionizing photon budget using updated models and calibrations. Building on previous work by Davies et al. [Davies2021] and Madau [Madau2017], we aim to assess whether star-forming galaxies can close the reionization photon budget at z\~{}7.
\section{Data and method}
Data and method:
To address this question, we adopt a literature-anchored approach, relying on published values for key parameters such as the cosmic star formation rate density (SFRD) from Madau \& Dickinson [Madau2014], ionizing photon production efficiency (xi\_ion), and escape fraction (f\_esc) calibrations. Specifically, we use xi\_ion = 10\^{}25.5 +/- 0.15 and f\_esc proxy calibrations based on O32/beta from LzLCS [Chisholm+22, Flury+22; Simmonds+24]. We perform a systematic reconciliation of these parameters to determine the required f\_esc for star-forming galaxies to close the reionization photon budget.
\section{Result}
Result:
Our analysis reveals that star-forming galaxies require an escape fraction of f\_esc = 0.105 (+0.106/-0.054) to reconcile the ionizing photon budget at z\~{}7, assuming a Madau-Dickinson SFRD and log xi\_ion = 25.5 +/- 0.15. This value is compared to the indirect-proxy-inferred f\_esc of 0.050 (+0.076/-0.030) from LzLCS O32/beta calibrations. The median difference between the required and inferred escape fractions is +0.048 dex-frac (16-84\%: -0.034 to +0.155), with 73\% of systematic Monte Carlo simulations showing a shortfall in ionizing photons.
\begin{figure}\centering\includegraphics[width=\columnwidth]{result.png}\caption{The reionization ionizing-photon-budget shortfall for star-forming galaxies is not robust to systematics at z\~{}7 (beta systematic corner)}\end{figure}
\section{Caveats}
Caveats:
Our result is subject to several limitations, primarily stemming from the reliance on published calibrations and assumptions about key parameters such as xi\_ion and clumping factor (C). The use of uncalibrated, automated measurements may introduce biases, and our analysis does not account for potential variations in these parameters across different galaxy populations. Additionally, the systematic uncertainties associated with the O32/beta proxy calibrations can significantly impact the inferred f\_esc values. Therefore, while our study provides a valuable reconciliation of the reionization photon budget, further investigation using direct observations and refined models is necessary to confirm these findings.
\section*{References}
\footnotesize
[Muoz2024] Muñoz et al. (2024). Reionization after JWST: a photon budget crisis?. bibcode 2024MNRAS.535L..37M \par [Park2022] Park et al. (2022). Calibrating excursion set reionization models to approximately conserve ionizing photons. bibcode 2022MNRAS.517..192P \par [Duncan2015] Duncan et al. (2015). Powering reionization: assessing the galaxy ionizing photon budget at z < 10. bibcode 2015MNRAS.451.2030D \par [Davies2021] Davies et al. (2021). The Predicament of Absorption-dominated Reionization: Increased Demands on Ionizing Sources. bibcode 2021ApJ...918L..35D \par [Madau2017] Madau (2017). Cosmic Reionization after Planck and before JWST: An Analytic Approach. bibcode 2017ApJ...851...50M

\end{document}
